\chapter{Perforated Eardrum}
\index{Perforated Eardrum}

\medskip
\noindent\textit{This chapter is for general education. Ear symptoms after injury, infection, or loud noise should be assessed when significant or persistent.}

\section*{What it is}
A perforated eardrum is a hole or tear in the thin membrane between the ear canal and middle ear. It can follow a middle-ear infection, a blow or object in the ear, a very loud noise, or a sudden pressure change such as diving or flying.

\section*{Symptoms}
Symptoms may include sudden ear pain, reduced hearing, ringing or buzzing, itching, dizziness, fever, or clear fluid, blood, or pus coming from the ear. Symptoms can be mild, and other conditions can cause similar hearing or pain problems.

\section*{Assessment and treatment}
A clinician can examine the ear and arrange hearing tests or other assessment if needed. Most perforations heal without an operation within several weeks to about 2 months. Antibiotics may be used for an infection. A persistent perforation or ongoing hearing loss may need an ear, nose, and throat assessment and sometimes repair surgery.

\section*{Self-care}
Keep the ear dry while it heals and do not swim unless advised. Do not put cotton buds, drops, oil, or other objects into the ear unless prescribed. Avoid forceful nose-blowing. Use pain medicine only as directed and ask a pharmacist or clinician if you have other conditions or medicines.

\section*{When to Seek Medical Care}
Seek urgent care for sudden hearing loss, worsening hearing, severe dizziness, facial weakness, severe pain, fever with swelling, or discharge after a head injury. Arrange medical review for persistent discharge, pain, or hearing loss, or if symptoms do not begin improving within a few weeks.

\section*{Sources and further reading}
\begin{itemize}
\item NHS. \textit{Perforated eardrum}. \url{https://www.nhs.uk/conditions/perforated-eardrum/}
\end{itemize}
